%%% FIELD prop-statement
For every compatible observed table \(p\), the following comparator reductions hold.  First, if \(\mathcal E=\{a,b\}\), then necessarily \(T=\mathcal E\), \(D_T(p)=u+v\), and \([L_T(p),U_T(p)]\) is literally the sharp Balke--Pearl interval for the binary exposure contrast \(\mathbb E[V_b-V_a]\).  Second, suppose \(T=\{a,b\}\) and \(U\neq\varnothing\).  Introduce a new label \(c^\dagger\) and form the three-category table
\[
 p^{{\rm m},z}_{ay}=p^z_{ay},\qquad
 p^{{\rm m},z}_{by}=p^z_{by},\qquad
 p^{{\rm m},z}_{c^\dagger y}=\sum_{x\in U}p^z_{xy}.
\]
Then \([L_T(p),U_T(p)]\) equals the eight-piece sharp interval computed from \(p^{\rm m}\), which is the setting and conclusion of Gabriel--Sachs--Sjolander Result 1: two trusted target levels and one merged wholly untrusted remainder.  Third, if \(\mathcal E=\{a,b,c\}\) and \(T=\mathcal E\), then the lower list consists of \(\ell_1,\ldots,\ell_8,C^-_{c,0},C^-_{c,1}\), the upper list consists of their outcome-reversed negatives, and these are exactly the ten lower and ten upper expressions in Gabriel--Sachs--Sjolander equations (3)--(4).  Finally, for arbitrary finite \(\mathcal E\) under full trust \(T=\mathcal E\), the image of the response-law model in pairs consisting of the joint law of \(V\) and the observed table \(p\) is the image of Song--Guo--Chan--Richardson model \(\mathcal M_1\), restricted to binary \(Z\) and \(Y\).  Hence its \(\theta_{ab}\)-projection is \([L_T(p),U_T(p)]\); by their equal-image theorem, exactly the same interval is the \(\theta_{ab}\)-projection under each of their models \(\mathcal M_2,\ldots,\mathcal M_5\).
%%% FIELD prop-proof
Because \(p^z\) is a conditional probability vector and \(0<\rho<1\),
\[
 p^z_{xy}\geq 0,
 \qquad
 \sum_{x\in\mathcal E}\sum_{y=0}^1p^z_{xy}=1,
 \qquad z\in\{0,1\}.
 \tag{1}
\]
Compatibility means that \(\mathcal Q_T(p)\neq\varnothing\).  The sharp-envelope theorem therefore applies and gives
\[
 \{\theta_{ab}(q):q\in\mathcal Q_T(p)\}
 =[L_T(p),U_T(p)],
 \quad
 L_T(p)=\max_{f\in\mathfrak F^-_T}f(p),
 \quad
 U_T(p)=\min_{f\in\mathfrak F^+_T}f(p).
 \tag{2}
\]

Suppose first that \(T=\{a,b\}\).  The index set \(T\setminus\{a,b\}\) is empty.  Thus the definition of the two-diagonal capacity gives
\[
 D_T(p)
 =\min\bigl(\{u+v\}\cup\varnothing\bigr)
 =u+v.
 \tag{3}
\]
The affine-list definition and (2) consequently give
\[
 L_T(p)=\max_{1\leq j\leq8}\ell_j(p),
 \qquad
 U_T(p)=\min_{1\leq j\leq8}\{-\ell_j(\bar p)\}.
 \tag{4}
\]
Every quantity on the right side of (4) is a function only of
\[
 \alpha_z=p^z_{a0},\quad \gamma_z=p^z_{a1},\quad
 \delta_z=p^z_{b0},\quad \beta_z=p^z_{b1},
 \qquad z\in\{0,1\}.
 \tag{5}
\]

If \(\mathcal E=\{a,b\}\), the observed exposure itself is binary.  Relabeling the two exposure values as \(0=a\) and \(1=b\), and using (5), term-by-term substitution in the definition of \(\ell_1,\ldots,\ell_8\) gives the eight Balke--Pearl lower expressions; substituting \(\bar p^z_{xy}=p^z_{x,1-y}\) gives their eight upper expressions as \(-\ell_j(\bar p)\).  The cited Balke--Pearl result states that the maximum and minimum of those respective lists are the sharp binary average-effect endpoints.  Combining that result with (2)--(5) proves the first reduction.  Notice that this is the only case in which the observed exposure alphabet is literally binary.

Now retain \(T=\{a,b\}\) but suppose \(U\neq\varnothing\), and define \(p^{\rm m}\) as in the statement.  Equation (1) shows that \(p^{{\rm m},z}\) is a probability vector.  It is also compatible with the three-category selected-trust model.  Indeed, choose any \(q\in\mathcal Q_T(p)\), put
\[
 X_z^{\rm m}=\begin{cases}
 X_z,&X_z\in\{a,b\},\\
 c^\dagger,&X_z\in U,
 \end{cases}
 \qquad V_a^{\rm m}=V_a,qquad V_b^{\rm m}=V_b,
 \tag{6}
\]
and, on the event \(X_z\in U\), set \(W^{\rm m}_{c^\dagger z}=W_{X_z z}\), assigning this binary slot arbitrarily off that event.  The events \(\{X_z=x\}\), \(x\in U\), are disjoint, so the compatible-response-law equations imply
\[
 \Pr(X_z^{\rm m}=c^\dagger,W^{\rm m}_{c^\dagger z}=y)
 =\sum_{x\in U}\Pr_q(X_z=x,W_{xz}=y)
 =\sum_{x\in U}p^z_{xy}
 =p^{{\rm m},z}_{c^\dagger y}.
 \tag{7}
\]
For \(x\in\{a,b\}\), (6) instead yields
\[
 \Pr(X_z^{\rm m}=x,V_x^{\rm m}=y)
 =\Pr_q(X_z=x,V_x=y)
 =p^z_{xy}
 =p^{{\rm m},z}_{xy}.
 \tag{8}
\]
Equations (7)--(8) prove compatibility of \(p^{\rm m}\) from the definition of \(\mathcal Q_{\{a,b\}}(p^{\rm m})\).

The four target-cell vectors in (5) are unchanged by merging.  Hence every term on the right side of (4), both for \(p\) and for its outcome reversal, is unchanged:
\[
 \ell_j(p^{\rm m})=\ell_j(p),
 \qquad
 \ell_j(\overline{p^{\rm m}})=\ell_j(\bar p),
 \qquad 1\leq j\leq8.
 \tag{9}
\]
Applying (2)--(4) to both compatible tables and using (9) proves equality of their sharp intervals.  Gabriel--Sachs--Sjolander Result 1 states that, with a binary instrument, two well-defining target exposure levels, and one merged potentially ill-defining level, these same first eight lower and upper expressions are sharp.  Thus this reduction is to their Result 1, rather than a claim that the original arbitrary-support observed exposure is binary.

For the ternary full-trust case, write \(\mathcal E=\{a,b,c\}\) and \(T=\mathcal E\).  There is exactly one trusted nuisance category.  From the definitions of \(\kappa\), \(h\), and \(C^-\),
\[
 \begin{aligned}
 C^-_{c,y}(p)
 &=\kappa-h_{cy}\\
 &=-2+\alpha_0+\alpha_1+\beta_0+\beta_1
      +p^0_{cy}+p^1_{c,1-y},
 \qquad y\in\{0,1\}.
 \end{aligned}
 \tag{10}
\]
Likewise, because \(C^+_{c,y}(p)=-C^-_{c,1-y}(\bar p)\),
\[
 C^+_{c,y}(p)
 =2-\gamma_0-\gamma_1-\delta_0-\delta_1
     -p^0_{cy}-p^1_{c,1-y},
 \qquad y\in\{0,1\}.
 \tag{11}
\]
The affine-list definition therefore supplies exactly eight baseline expressions plus the two expressions in (10) for the lower endpoint, and eight outcome-reversed baseline expressions plus the two expressions in (11) for the upper endpoint.  Under the relabeling \(x=a\), \(x'=b\), and \(x''=c\), direct substitution of \(p_{xy\cdot z}=p^z_{xy}\) shows that (10) and (11) are respectively the final two lines of Gabriel--Sachs--Sjolander equations (3) and (4), while the definitions of \(\ell_1,\ldots,\ell_8\) give their first eight lines.  Their cited ternary result establishes that those ten-term maxima and minima are sharp.  Together with (2), this proves the third reduction.

Finally suppose \(T=\mathcal E\).  Then \(U=\varnothing\), and the response space and observation equations reduce to
\[
 \mathcal R_T=\mathcal E^2\times\{0,1\}^{\mathcal E},
 \qquad R=(X_0,X_1,V),
 \qquad X=X_Z,
 \qquad Y=V_X.
 \tag{12}
\]
For any \(q\in\mathcal Q_T(p)\), draw \(Z\) independently of \(R\), as required by random assignment and the mixed-trust observation map.  Equations (12) and the definition of compatibility then give
\[
 \Pr_q(X_z=x,V_x=y)=p^z_{xy},
 \qquad x\in\mathcal E,quad y,z\in\{0,1\}.
 \tag{13}
\]
Thus \((X_0,X_1,V)\), together with the independent instrument, is a law in Song--Guo--Chan--Richardson model \(\mathcal M_1\): it has exposure responses for both instrument values, a single outcome response \(V_x\) for every exposure value, individual-level exclusion, consistency, and random assignment.

Conversely, take a law in their \(\mathcal M_1\), restricted to binary \(Z,Y\) and exposure alphabet \(\mathcal E\).  Its marginal law of the complete response vector \((X(0),X(1),(Y(x):x\in\mathcal E))\) defines a \(q\) on the space in (12).  Individual-level exclusion and consistency identify \(Y=Y(X)=V_X\), while random assignment and \(0<\rho<1\) imply
\[
 \Pr_q(X_z=x,V_x=y)
 =\Pr(X=x,Y=y\mid Z=z)
 =p^z_{xy}.
 \tag{14}
\]
Hence \(q\in\mathcal Q_T(p)\).  Equations (13)--(14) prove equality, at every fixed \(p\), between the attainable joint laws of \(V\) in the full-trust model here and the corresponding fiber of the image of \(\mathcal M_1\).

The cited equal-image theorem of Song--Guo--Chan--Richardson states that the image pairs consisting of the joint potential-outcome law and the observed conditional law are identical for \(\mathcal M_1,\ldots,\mathcal M_5\).  Let \(\phi\) denote their image map and, for \(i\in\{1,\ldots,5\}\), define its fiber at the observed table by
\[
 \Phi_i(p)=\{\mu:(\mu,p)\in\phi(\mathcal M_i)\},
 \tag{15}
\]
where \(\mu\) is a probability mass function on \(\{0,1\}^{\mathcal E}\).  Equality of the pair images gives, by taking the section with second coordinate \(p\),
\[
 \Phi_1(p)=\Phi_2(p)=\cdots=\Phi_5(p).
 \tag{16}
\]
Define the linear functional
\[
 \tau_{ab}(\mu)=\sum_{v\in\{0,1\}^{\mathcal E}}(v_b-v_a)\mu(v).
 \tag{17}
\]
This functional depends only on the joint potential-outcome law.  Applying \(\tau_{ab}\) to the equal fibers in (16), and applying the sharp-envelope theorem to the \(\mathcal M_1\) fiber via (12)--(14), gives, for every \(i\in\{1,\ldots,5\}\),
\[
 \{\tau_{ab}(\mu):\mu\in\Phi_i(p)\}
 =[L_T(p),U_T(p)]
 \tag{18}
\]
because \(\tau_{ab}(\operatorname{Law}_q(V))=\mathbb E_q[V_b-V_a]=\theta_{ab}(q)\).  This proves the last claim.
%%% FIELD cited-balke-statement
For a binary instrument, binary exposure, and binary outcome under random assignment, individual-level exclusion, and consistency, Balke and Pearl's displayed average-causal-effect bounds are sharp and consist of a lower endpoint equal to the maximum of eight affine observed-cell expressions and an upper endpoint equal to the minimum of eight affine observed-cell expressions.
%%% FIELD cited-gabriel-result-one-statement
For a binary instrument, two well-defining exposure levels used in the causal contrast, and one merged potentially ill-defining exposure level, Gabriel, Sachs, and Sjolander Result 1 gives sharp bounds that are symbolically identical to the first eight lower and first eight upper expressions of their equations (3)--(4), namely the classical binary-instrument bounds.
%%% FIELD cited-gabriel-ternary-statement
For a binary instrument, binary outcome, and exactly three well-defining exposure levels, Gabriel, Sachs, and Sjolander equations (3)--(4) give the sharp contrast endpoints as the maximum of ten affine lower expressions and the minimum of ten affine upper expressions: the eight classical expressions plus two expressions involving the third exposure level on each side.
%%% FIELD cited-song-equal-image-statement
Writing \(\mathcal M_1,\ldots,\mathcal M_5\) for the five finite categorical instrumental-variable models in Song, Guo, Chan, and Richardson Table 2, and using their equation (13) map from a model law to the pair consisting of the joint outcome-potential law and the observed conditional law, Theorem 1 states that all five models have the same image, characterized by their inequalities (8).
